\documentclass[11pt]{article}

\usepackage[margin=1in]{geometry}
\usepackage{amsmath,amssymb,amsthm}
\usepackage{booktabs}
\usepackage[hidelinks]{hyperref}
\hypersetup{
  pdftitle={The Width-Four q-Fibonomial Coefficients Are Unimodal},
  pdfauthor={hainzhao}
}

\newtheorem{theorem}{Theorem}
\newtheorem{lemma}[theorem]{Lemma}
\newcommand{\qint}[1]{[#1]_q}
\newcommand{\qfib}[2]{\left[\!\begin{matrix}#1\\#2\end{matrix}\!\right]_{\!\mathcal F}}
\newcommand{\archivedoi}{10.5281/zenodo.21826970}

\title{The Width-Four \(q\)-Fibonomial Coefficients Are Unimodal}
\author{hainzhao}
\date{August 2026\\[4pt]
\small Proof and source archive:
\href{https://doi.org/\archivedoi}{doi:\archivedoi}}

\begin{document}
\maketitle

\begin{abstract}
Bergeron, Ceballos, and K\"ustner conjectured that every
\(q\)-Fibonomial coefficient is unimodal.  Recent work proves the cases of
width at most three.  We prove the next infinite family: for every integer
\(m\geq1\), the polynomial \(\qfib{m+4}{4}\) is unimodal.  After taking a
first coefficient difference, the problem reduces below the midpoint to four
translates of the restricted partition function for parts \(1,2,3\).  Its
period-six quasipolynomial formula yields elementary quadratic lower bounds
for all \(m\geq8\); the seven remaining cases are checked exactly.
\end{abstract}

\section{Introduction}

Let \(F_0=0,F_1=1\), and \(F_{j+2}=F_{j+1}+F_j\).  For a positive integer
\(r\), write
\[
  \qint r=1+q+\cdots+q^{r-1},\qquad
  [F_r]^!_q=\prod_{j=1}^r [F_j]_q.
\]
Building on the unweighted path--domino model of Sagan and Savage
\cite{SaganSavage}, Bergeron, Ceballos, and K\"ustner define
\[
  \qfib{m+n}{n}
  =\frac{[F_{m+n}]^!_q}{[F_m]^!_q[F_n]^!_q}.
  \tag{1}
\]
Their path--domino model proves that (1) is a polynomial with nonnegative
integer coefficients \cite[Theorem~2.4]{BCK}; they conjecture that all these
polynomials are unimodal \cite[Conjecture~2.5]{BCK}.  Connelly, Ito,
Martinez, Shevchenko, and Yang prove the conjecture for \(n\leq3\)
\cite[Theorem~1.2]{CIMSY} and report finite computations for larger parameters.

This note proves the width-four case.  The proof is algebraic and uses the
published path--domino theorem only to import polynomiality and nonnegativity.
The May 2026 preprint \cite{CIMSY} explicitly leaves the cases \(n\geq4\)
open and notes that its proposed product criterion does not directly fit
their factorization.  Our literature search through August 6, 2026 found no
later proof of width four; this is a bounded search statement, not an
exhaustive priority claim.

\begin{theorem}\label{thm:main}
For every integer \(m\geq1\), the \(q\)-Fibonomial coefficient
\(\qfib{m+4}{4}\) is unimodal.
\end{theorem}

Here a polynomial \(C(q)=\sum_{j=0}^D c_jq^j\) is \emph{unimodal} if its
coefficient sequence first weakly increases and then weakly decreases.

\section{The midpoint difference formula}

Fix \(m\geq1\), and put
\[
  a=F_{m+1},\qquad b=F_{m+2}.
\]
The next two Fibonacci numbers are \(a+b\) and \(a+2b\).  Since
\([F_1]_q=[F_2]_q=1\), \([F_3]_q=[2]_q\), and
\([F_4]_q=[3]_q\), equation (1) becomes
\[
 W_m(q):=\qfib{m+4}{4}
 =\frac{[a]_q[b]_q[a+b]_q[a+2b]_q}{[2]_q[3]_q}.
 \tag{2}
\]
Its degree is
\[
 D=3a+4b-7. \tag{3}
\]
Indeed, \([r]_{q^{-1}}=q^{1-r}[r]_q\), so (2) gives
\(W_m(q^{-1})=q^{-D}W_m(q)\).  The path--domino theorem makes \(W_m\) a
polynomial, and its constant term is one; consequently it has degree \(D\)
and is symmetric.

For an integer \(t\), define
\[
 p(t)=[q^t]\frac{1}{(1-q)(1-q^2)(1-q^3)}.
\]

\begin{lemma}\label{lem:p}
One has
\[
 p(t)=
 \begin{cases}
 \displaystyle\left\lfloor\frac{t^2+6t+12}{12}\right\rfloor,&t\geq0,\\[6pt]
 0,&t<0.
 \end{cases} \tag{4}
\]
In particular, for \(t\geq0\),
\[
 \frac{t^2+6t+1}{12}\leq p(t)
 \leq\frac{t^2+6t+12}{12}. \tag{5}
\]
\end{lemma}

\begin{proof}
For \(t\geq0\), the generating function shows that \(p(t)\) counts the
solutions of \(x+2y+3z=t\) in nonnegative integers.  Thus
\[
 p(t)=\sum_{z=0}^{\lfloor t/3\rfloor}
 \left(\left\lfloor\frac{t-3z}{2}\right\rfloor+1\right).
\]
Writing \(t=6u+r\), \(0\leq r<6\), and summing separately over even and
odd \(z\) gives (4) in each of the six residues.  The upper bound in (5) is
immediate.  Since the numerator in (4) is an integer, rounding down loses at
most \(11/12\), which gives the lower bound.
\end{proof}

Write \(W_m(q)=\sum_{t=0}^Dc_tq^t\), with \(c_{-1}=0\), and let
\(T=\lfloor D/2\rfloor\).  Symmetry reduces unimodality to showing
\(c_t-c_{t-1}\geq0\) for \(0\leq t\leq T\).  From (2),
\[
 (1-q)W_m(q)=
 \frac{(1-q^a)(1-q^b)(1-q^{a+b})(1-q^{a+2b})}
 {(1-q)(1-q^2)(1-q^3)}. \tag{6}
\]

\begin{lemma}\label{lem:difference}
For \(0\leq t\leq T\),
\[
 c_t-c_{t-1}
 =p(t)-p(t-a)-p(t-b)+p(t-(2a+b)). \tag{7}
\]
\end{lemma}

\begin{proof}
Expand the numerator of (6) by inclusion--exclusion.  The pair shift
\(a+b\) cancels the singleton shift \(a+b\), while the pair shift
\(b+(a+b)=a+2b\) cancels the singleton shift \(a+2b\).  The only other pair
shift that can occur at or below \(T\) is \(a+(a+b)=2a+b\).  Indeed, every
remaining pair shift is at least \(2a+2b>T\), as follows directly from (3).
Every triple shift is also at least \(2a+2b>T\).  The four surviving terms
give (7).
\end{proof}

\section{Uniform positivity}

We now prove the required coefficient inequalities for \(m\geq8\).  In this
range, \(a\geq F_9=34\).  Let the right side of (7) be \(g(t)\).

If \(0\leq t<a\), then \(g(t)=p(t)>0\).  If \(a\leq t<b\), then
\(g(t)=p(t)-p(t-a)\geq0\), since the counting function \(p\) is
nondecreasing.

Suppose next that \(b\leq t<2a+b\).  The last term of (7) vanishes.  Applying
the lower estimate in (5) to the positive term and the upper estimate to the
two negative terms gives
\[
 12g(t)\geq K(t), \tag{8}
\]
where
\[
 K(t)=-t^2+(2a+2b-6)t-a^2-b^2+6a+6b-23.
\]
The quadratic \(K\) is concave, so its minimum on the integer interval
\([b,2a+b-1]\) occurs at an endpoint.  We have
\[
 K(b)=a(2b-a+6)-23,
\]
and
\[
 K(2a+b-1)>K(2a+b)=a(2b-a-6)-23, \tag{9}
\]
where the first difference is exactly \(2a+5\).  Since \(b>a\) and
\(a\geq34\), the displayed expressions for \(K(b)\) and \(K(2a+b)\) are
positive.  Thus \(g(t)>0\) throughout this interval.

Finally, suppose \(2a+b\leq t\leq T\).  Applying the appropriate side of
(5) to all four terms of (7) yields
\[
 12g(t)\geq -2at+3a^2+4ab-6a-22. \tag{10}
\]
The right side decreases with \(t\).  By (3),
\[
 t\leq T\leq\frac{3a+4b-7}{2},
\]
and substitution in (10) gives
\[
 12g(t)\geq a-22>0. \tag{11}
\]
This proves all midpoint differences for \(m\geq8\).

\section{The seven finite cases}

For \(1\leq m\leq7\), we evaluate the exact integer formula (7) over the
entire interval \(0\leq t\leq T\).  The resulting minima are shown in
Table~\ref{tab:small}.

\begin{center}
\refstepcounter{table}\label{tab:small}
\centering
\begin{tabular}{cccr}
\toprule
\(m\)&\((a,b)\)&\(T\)&\(\min_{0\leq t\leq T}g(t)\)\\
\midrule
1&(1,2)&2&0\\
2&(2,3)&5&0\\
3&(3,5)&11&0\\
4&(5,8)&20&0\\
5&(8,13)&34&1\\
6&(13,21)&58&1\\
7&(21,34)&96&1\\
\bottomrule
\end{tabular}

\smallskip
\small Table~\thetable: Exact finite cases not covered by the uniform estimate.
\end{center}

Every listed minimum is nonnegative.  Together with the preceding section,
Lemma~\ref{lem:difference} shows that the coefficients weakly increase to the
midpoint for every \(m\geq1\).  Symmetry then proves
Theorem~\ref{thm:main}.

\section{Reproducibility and claim boundary}

The finite calculation uses only integer arithmetic.  From the repository
root, it is replayed by
\begin{verbatim}
python3 proof/qfib_width4_unimodality_proof.py
\end{verbatim}
The script performs three checks: it compares (4) with direct enumeration for
\(0\leq t\leq300\); it verifies every entry of Table~\ref{tab:small}; and,
as an independent consistency check, it constructs the exact quotient (2)
for \(1\leq m\leq24\) and compares all midpoint differences with (7).  The
last two ranges beyond \(m=7\) are checks, not assumptions in the proof.

The theorem proves only the family \(\qfib{m+4}{4}\), \(m\geq1\).  It does
not prove the full two-parameter conjecture, a combinatorial chain
decomposition, log-concavity, or any width \(n\geq5\).

\paragraph{AI assistance.}
The author used AI assistance in preparing this draft and its verification
code.

\bibliographystyle{plain}
\bibliography{references}

\end{document}
